% Final integrated draft produced from main.tex, draft.tex, and current result files.

\documentclass{vgtc}

\graphicspath{{figures/}{pictures/}{images/}{./}}

\usepackage{times}
\renewcommand*\ttdefault{txtt}
\usepackage{mathptmx}
\usepackage{amsmath}
\usepackage{amssymb}
\usepackage{array}
\usepackage{booktabs}
\usepackage{multirow}
\usepackage{graphicx}
\usepackage{xcolor}
\usepackage{url}
\usepackage{enumitem}
\usepackage{xspace}

\newcommand{\R}{\mathbb{R}}
\newcommand{\Domain}{\mathcal{M}}
\newcommand{\Range}{\mathcal{R}}
\newcommand{\Sheets}{\mathcal{S}}
\newcommand{\Sheet}{S}
\newcommand{\Score}{C}
\newcommand{\placeholder}[2]{%
  \begin{center}
  \fbox{%
    \begin{minipage}[c][#1][c]{0.92\linewidth}
      \centering
      #2
    \end{minipage}}
  \end{center}
}
\newcommand{\wideplaceholder}[2]{%
  \begin{center}
  \fbox{%
    \begin{minipage}[c][#1][c]{0.96\textwidth}
      \centering
      #2
    \end{minipage}}
  \end{center}
}
\newcommand{\myparagraph}[1]{\vspace{1mm} \noindent \textbf{#1}}
\newcommand{\ie}{i.e.,\xspace}
\newcommand{\etal}{et al.\xspace}
\newcommand{\Mohit}[1]{\textcolor{blue}{#1}\xspace}
\newcommand{\Talha}[1]{\textcolor{magenta}{Talha: #1}\xspace}
\newcommand{\Nanna}[1]{\textcolor{purple}{#1}\xspace}
\newcommand{\Ingrid}[1]{\textcolor{red}{#1}\xspace}
\newcommand{\Petar}[1]{\textcolor{cyan}{#1}\xspace}

\onlineid{0}
\vgtccategory{Research}
\vgtcinsertpkg

\title{Topological Tracking of Bivariate Features Using Reeb Spaces}

\author{Mohit Sharma\thanks{e-mail: mohit.sharma@liu.se}\\ %
        \scriptsize Link\"oping university %
\and Petar Hristov\thanks{e-mail: petar.georgiev.hristov@liu.se}\\ %
        \scriptsize Link\"oping university %
\and Nanna Holmgaard List \thanks{e-mail: nalist@kth.se}\\ %
        \scriptsize KTH Royal Institute of Technology %
\and Talha Bin Masood\thanks{e-mail: talha.bin.masood@liu.se}\\ %
        \scriptsize Link\"oping university %
\and Ingrid Hotz\thanks{e-mail: ingrid.hotz@liu.se}\\ %
        \scriptsize Link\"oping university %
}

\teaser{
  \centering
  \wideplaceholder{1.45in}{[Teaser figure: one molecular sequence window with the range-event graph, a linked Sankey interval, selected sheet footprints, and corresponding fiber-surface images. The visual message should be that a Reeb-space sheet can be followed as a temporal feature and inspected simultaneously in range and domain context.]}
  \caption{Our workflow tracks prominent Reeb-space sheets in time-varying bivariate fields and links temporal events to sheet and fiber-surface evidence.}
  \label{fig:teaser-final}
}

\abstract{
Time-varying bivariate fields occur in molecular electronic-structure analysis, synthetic multifield benchmarks, and many other simulations where the relationship between two scalar quantities changes over time. Existing continuous-scatterplot summaries describe how the distribution of bivariate values evolves, but they do not directly provide feature tracks for individual topological structures. We present a visual analytics workflow for tracking prominent sheets of the Reeb space of a bivariate field. For each timestep, we compute a Reeb space, rank sheets by range-space area, and retain the top $N$ sheets as trackable bivariate features. We compute candidate correspondences between adjacent timesteps using two complementary notions of continuity: range-space similarity of sheet footprints and domain-space overlap of regular domain vertices. The resulting temporal graph is shown in an interactive Sankey view linked to range-space sheet renderings, fiber-surface renderings, event diagnostics, continuing-feature summaries, stride controls, and domain/range disagreement views. We evaluate the workflow on two real molecular datasets, methyl vinyl ketone and stilbene, and one controlled torus dataset. Range-space sheet tracking recovers the known MVK transition window, identifies a robust late stilbene reconfiguration, and localizes the controlled torus transition. Domain-space rankings provide a separate view of domain-vertex overlap changes, while target disagreement highlights intervals where range and domain criteria select different continuations. These results show that Reeb-space sheets are useful as trackable topological features for visual analysis of time-varying bivariate data.
}

\keywords{Bivariate fields, Reeb space, feature tracking, fiber surfaces, time-varying data, molecular visualization.}

\begin{document}

\firstsection{Introduction}
\maketitle

Many scientific simulations produce multiple scalar fields on a common spatial domain. In molecular electronic-structure analysis, for example, an electronic transition can be represented by hole and particle natural transition orbital fields~\cite{Martin2003NTO}. These two fields indicate where charge is depleted and accumulated during the transition. When molecular geometry changes over time, the analysis problem becomes time-varying and bivariate: the analyst wants to understand how the joint behavior of two scalar fields evolves, not only how either field evolves in isolation.

Continuous scatterplots (CSPs) are a natural way to summarize bivariate fields because they show the distribution of values in the two-dimensional range. Prior work used CSP descriptors to study temporal changes in molecular electronic-structure data~\cite{Sharma2025}. That approach gives useful timestep-level signatures, and domain-segmentation-driven CSP peeling can expose meaningful bivariate components~\cite{Sharma2021}. However, a CSP descriptor does not directly answer feature-tracking questions such as: Which topological feature persists? Which feature disappears? Which interval contains a reconfiguration? Which domain vertices continue to participate in the same bivariate behavior?

We address these questions using Reeb spaces. For a bivariate field $F=(f,g):\Domain\rightarrow\R^2$, the Reeb space identifies points that lie in the same connected component of a fiber. Its two-dimensional cells, called sheets, combine range-space geometry with domain connectedness. A sheet can therefore be interpreted as a coherent bivariate feature. Tracking sheets across time gives a feature-level temporal representation that complements CSP descriptors.

This paper presents a practical visual analytics workflow for Reeb-space sheet tracking. At each timestep, we compute a Reeb space and retain the top $N$ sheets by range-space area. We compare sheets between adjacent timesteps using range-space shape similarity and domain-space vertex overlap. The resulting temporal graph is explored through a Sankey-based view, linked sheet images, linked fiber-surface images, and diagnostic plots. The tool does not force a one-to-one matching, because Reeb-space sheets can split, merge, appear, disappear, and reconfigure. Instead, it exposes candidate correspondences and lets the analyst compare how different continuity notions behave.

The key idea is that range and domain continuity answer different questions. Range-space similarity asks whether the same relationship between the two scalar fields persists. Domain-space overlap asks whether the same domain vertices continue to participate in the sheet. These questions need not have the same answer. We therefore report range events, domain events, and domain/range target disagreement separately. The disagreement view is not a replacement for the event ranking; it is a mechanism for selecting concrete examples where the two continuity notions confidently choose different targets.

Our contributions are:
\begin{itemize}[leftmargin=*]
  \item A Reeb-space sheet tracking workflow for time-varying bivariate fields, where prominent sheets serve as trackable topological features.
  \item Complementary correspondence measures for range-space sheet geometry and domain-space vertex overlap, with diagnostic summaries that keep the two interpretations separate.
  \item A linked visual interface combining Sankey-based temporal summaries, event plots, continuing-feature lists, stride controls, range-space sheet images, fiber-surface views, and domain/range disagreement highlights.
  \item Case studies on two real molecular datasets, MVK and stilbene, and one torus dataset showing that range-space sheet events recover meaningful reconfiguration intervals, while domain-space events expose a separate view of domain-vertex overlap changes.
\end{itemize}

\begin{figure*}[t]
  \wideplaceholder{1.55in}{[Figure content: show the domain $\Domain$ as a tetrahedral mesh, a fiber and fiber surface, the bivariate range with the arrangement induced by singular edges, the resulting Reeb-space sheets, and the corresponding domain regions. Use this as the conceptual figure that explains why sheets can be treated as features.]}
  \caption{Topological objects used in the paper. A bivariate field maps a spatial domain to a two-dimensional range. The Reeb space groups connected fiber components; its sheets represent coherent bivariate features with both range-space footprints and associated domain vertices.}
  \label{fig:topological-objects}
\end{figure*}

\section{Related Work}

\myparagraph{Reeb spaces and multivariate topology.}
The Reeb space generalizes the Reeb graph from scalar functions to multivariate maps~\cite{Edelsbrunner2008}. In the bivariate case, it describes how connected components of fibers change over the range and can be represented as a two-dimensional sheet complex connected along singular structures~\cite{Edelsbrunner2004b}. Algorithms for computing bivariate Reeb spaces include Jacobi fiber surface methods and arrange-and-traverse approaches~\cite{Tierny2017,Hristov2026SingularAA}. Approximate alternatives, including the Joint Contour Net and Mapper, have been used when exact Reeb-space computation is expensive or when robustness is more important than full cell-level detail~\cite{Carr2014,Dey2016,Brown2021}. Reeb spaces have also been used as descriptors for comparing multifields~\cite{Ramamurthi2022a,Agarwal2021,Ramamurthi2024}. Our work uses computed Reeb-space sheets not only as static descriptors, but as trackable features in a time-varying visual analytics workflow.

\myparagraph{Feature tracking.}
Topology-based feature tracking often follows one of two strategies: extract features independently at each timestep and match them afterward, or incorporate time into the domain and extract spatio-temporal structures directly~\cite{Post2003,Weber2011}. Many methods track critical points, merge-tree components, contour-tree features, vortices, or Morse-Smale structures using spatial overlap, persistence information, tree edit distances, or optimization-based correspondence models~\cite{Samtaney1994,silver1998tracking,Reininghaus2011,Sridharamurthy2020,Wetzels2022a,Wetzels2023}. Region-based tracking frequently uses spatial overlap and geometric attributes to identify continuations, splits, and merges. Our setting differs because the tracked objects are Reeb-space sheets of bivariate fields. They have a range-space footprint and associated domain vertices, and either view may be the more relevant notion of continuity depending on the analysis question.

\myparagraph{Bivariate molecular visualization.}
Fiber surfaces and CSPs are widely used to inspect bivariate data. Fiber surfaces provide detailed spatial context for selected range values, while CSPs summarize the distribution of bivariate values over the domain. Recent molecular visualization work used CSP image moments to study temporal transitions in excited-state fields~\cite{Sharma2025}. Our method is complementary: we retain the range-space perspective, but add Reeb-space sheets as feature units that can be connected across timesteps and inspected through linked sheet and fiber-surface views.

\section{Background}

\myparagraph{Bivariate fields.}
Let $\Domain$ be a tetrahedral mesh or, more generally, a triangulated three-dimensional domain. A bivariate field is a piecewise-linear map $F=(f,g):\Domain\rightarrow\R^2$, where $f$ and $g$ are scalar fields defined on the same vertices and interpolated over simplices. We refer to $\Domain$ as the domain and to $\R^2$ as the range. In the real molecular datasets, $f$ and $g$ are the configured orbital fields \texttt{orb00} and \texttt{orb01}.

\myparagraph{Fibers and fiber surfaces.}
For a range value $y=(a,b)$, the fiber $F^{-1}(y)$ is the set of domain points mapped to $y$. In a generic bivariate map over a three-dimensional domain, fibers are typically one-dimensional curves. For a region $P\subset\R^2$, the preimage $F^{-1}(P)$ sweeps out a fiber surface. In our viewer, fiber-surface images provide the spatial context needed to interpret selected range-space sheets and sheet correspondences.

\myparagraph{Reeb spaces and sheets.}
The Reeb space of $F$ is the quotient $\Domain/\sim$, where $x\sim x'$ when $F(x)=F(x')$ and $x$ and $x'$ lie in the same connected component of the fiber $F^{-1}(F(x))$. For a bivariate map over a three-dimensional domain, the Reeb space is represented as a two-dimensional sheet complex. Each sheet has a polygonal footprint in the bivariate range and is associated with a subset of regular domain vertices. We treat these sheets as the feature units to be tracked.

\section{Workflow Overview}

\begin{figure*}[t]
  \wideplaceholder{1.55in}{[Figure content: pipeline overview. Show the stage sequence from VTU inputs to Reeb-space files, RSI JSON, cached sheet geometry, range-shape matches, domain-overlap matches, sheet/fiber-surface images, analysis diagnostics, and the unified viewer.]}
  \caption{Pipeline overview. The method computes Reeb-space sheets per timestep, tracks prominent sheets with range and domain correspondence measures, renders the visual evidence, and presents the results in a linked visual analytics interface.}
  \label{fig:pipeline-overview}
\end{figure*}

The workflow starts from a time-ordered sequence of volumetric datasets. Each timestep contains the same two scalar arrays, giving a bivariate field $F_i=(f_i,g_i)$. The implementation is organized as a set of reproducible stages. Stage 1 runs the Reeb-space executable on each \texttt{.vtu} input and produces \texttt{.rs} and \texttt{.rsi} files. The \texttt{.rs} file stores the Reeb-space geometry and sheet structure used for geometry export and rendering. The \texttt{.rsi} file stores sheet information used by the Python stages, including sheet identifiers, areas, ranks, and regular domain vertices. Timesteps without usable Reeb-space sheet data are excluded from the tracking graph.

Stage 2 converts the binary \texttt{.rsi} files into compact per-timestep JSON summaries. Let $\Sheets_i^\mathrm{all}$ be all sheets at timestep $i$. We rank sheets by range-space area and retain the top $N$ sheets:
$$
\Sheets_i = \operatorname{TopN}(\Sheets_i^\mathrm{all}, N).
$$
All experiments in this paper use $N=20$. This choice keeps the temporal graph readable while preserving the dominant bivariate features. Sheets with no stored regular vertices are retained because their range-space footprint and area can still be meaningful; their domain-overlap evidence is zero by construction.

Stage 3A exports cached sheet geometry and computes range-space shape matches between selected sheets in adjacent timesteps. Stage 3B computes domain-vertex overlaps between the same selected sheet sets and attaches the available range metrics to the domain-overlap links. These stages produce the candidate correspondence graph used by the viewer and by the analysis diagnostics. Stage 4A renders sheet footprint images using common global range-space bounds and fixed image size, preventing misleading comparisons caused by per-image rescaling. Stage 4B renders top-sheet fiber-surface images, which connect the abstract range-space sheet to its spatial embedding in the domain. Stage 5 writes the self-contained browser viewer. A separate analysis stage consumes the generated viewer data and exports event, lifetime, sensitivity, metric-agreement, and domain/range disagreement summaries; it does not recompute Reeb spaces, sheet geometry, or images.

\section{Sheet Correspondence}

For each adjacent timestep pair $(i,i+1)$, every source sheet $\Sheet\in\Sheets_i$ is compared with every target sheet $\Sheet'\in\Sheets_{i+1}$. We store all candidate correspondences rather than solving a global one-to-one assignment. This design is intentional: sheet evolution can be many-to-many, and retaining the candidate graph supports visual inspection of split, merge, and ambiguous continuation hypotheses.

\myparagraph{Range-space similarity.}
The primary range-space metric is sheet footprint intersection-over-union. Let $M_\Sheet$ and $M_{\Sheet'}$ be raster masks of two sheet footprints under common global range bounds. The range score is
$$
q_R(\Sheet,\Sheet') =
\frac{|M_\Sheet \cap M_{\Sheet'}|}{|M_\Sheet \cup M_{\Sheet'}|}.
$$
This is the default and current range score in the analysis. The implementation also computes area ratio, bounding-box IoU, and centroid similarity. We do not use those auxiliary scores as primary result scores because mask IoU already combines the most important range-footprint evidence: overlap, relative area, position, and detailed shape under a common coordinate frame. Area ratio ignores position, centroid similarity ignores shape, and bounding-box IoU is coarser than the actual footprint. We keep them as diagnostics and optional viewer metrics; their definitions and agreement statistics are given in the supplement.

\myparagraph{Domain-space similarity.}
Let $V_\Sheet$ and $V_{\Sheet'}$ be the regular domain vertices associated with two sheets. We compute the raw overlap $O(\Sheet,\Sheet')=|V_\Sheet\cap V_{\Sheet'}|$ and two directional overlap ratios:
$$
d_\mathrm{src}(\Sheet,\Sheet') =
\begin{cases}
O(\Sheet,\Sheet')/|V_\Sheet|, & |V_\Sheet|>0,\\
0, & |V_\Sheet|=0,
\end{cases}
\qquad
d_\mathrm{tgt}(\Sheet,\Sheet') =
\begin{cases}
O(\Sheet,\Sheet')/|V_{\Sheet'}|, & |V_{\Sheet'}|>0,\\
0, & |V_{\Sheet'}|=0.
\end{cases}
$$
A natural symmetric alternative is vertex-set Jaccard overlap,
$$
q_\mathrm{jac}(\Sheet,\Sheet') =
\frac{|V_\Sheet\cap V_{\Sheet'}|}{|V_\Sheet\cup V_{\Sheet'}|}.
$$
We keep this definition as an important negative design choice. Jaccard overlap is useful when the question is whether two vertex sets are nearly identical, but that is not the main domain question in our tracking problem. If a sheet grows, shrinks, is partly absorbed, or contributes to a merge, the same source vertices may continue into the target while the union becomes much larger. Jaccard would then score a meaningful containment continuation as weak. The directional ratios separate the two interpretations: $d_\mathrm{src}$ asks how much of the source survives, and $d_\mathrm{tgt}$ asks how much of the target is inherited from the source.
The domain event and disagreement views use $q_D=d_\mathrm{max}=\max(d_\mathrm{src},d_\mathrm{tgt})$. This ratio lies in $[0,1]$ and treats containment as strong evidence of domain continuity, which is useful when a feature grows, shrinks, merges, or is partially absorbed. UI labels may display the same value as a percentage, but all formulas use the ratio.

\section{Visual Analytics Interface}

\begin{figure*}[t]
  \wideplaceholder{1.75in}{[Figure content: unified viewer screenshot. Include two Sankey panels comparing shape-overlap and domain-overlap views, a selected interval in the event graph, highlighted source and target nodes, side-by-side sheet images, side-by-side fiber-surface images, and the linked zoom/image detail panel.]}
  \caption{The interface links temporal Reeb-sheet correspondences to range-space and domain-space evidence. Analysts can compare shape-overlap and domain-overlap panels, select event intervals and continuing features, inspect threshold sensitivity, highlight domain/range disagreement examples, and examine sheets and fiber surfaces for individual nodes and links.}
  \label{fig:viewer}
\end{figure*}

The main visualization is a Sankey-style temporal graph. Timesteps are placed from left to right, selected sheets are shown as nodes, and candidate correspondences are shown as links. Link thickness and opacity encode the active correspondence score after normalization within the selected metric mode. Analysts can filter links by threshold, display only the strongest outgoing link per source sheet, change the visible timestep range and stride, adjust the number of top sheets, and reorder sheets using weighted crossing-reduced ordering, sheet area, or domain vertex count.

The interface is designed for comparison. Multiple Sankey panels can be shown side by side, each with its own data mode, metric, threshold, and visual scaling. A typical configuration places shape-overlap correspondences next to domain-overlap correspondences. Differences between panels reveal whether the same bivariate range feature persists, whether the same domain vertices persist, or whether the two notions diverge.

The analysis panel turns the graph into an exploratory workflow. The interval tab ranks timestep pairs with weak continuations or possible split/merge ambiguity. The continuing-feature tab follows best outgoing links above the selected threshold. The sensitivity tab repeats these summaries over threshold values. The domain/range complementarity tab highlights source sheets for which range and domain criteria confidently choose different targets. Node and link selections open linked detail views: a node shows sheet metadata, the rendered sheet footprint, and available fiber-surface images; a link shows source and target sheets side by side. Image views support linked pan and zoom so that corresponding sheet or fiber-surface structures can be compared at the same scale. The supplement gives the exhaustive control-level description.

\section{Diagnostics}

\myparagraph{Event intervals.}
For a selected score $q$, each source sheet has a best outgoing continuation and each target sheet has a best incoming predecessor. At threshold $\theta$, we count weak outgoing sheets, weak incoming sheets, possible splits, and possible merges. The event score for interval $i\rightarrow i+1$ is
$$
E_i(\theta)=
W^\mathrm{out}_i + W^\mathrm{in}_i
+0.5(P^\mathrm{split}_i+P^\mathrm{merge}_i)
+(1-\overline{B^\mathrm{out}_i})|\Sheets_i|.
$$
Large values indicate intervals where many sheets lack strong continuations, many targets lack strong predecessors, branching is ambiguous, or the average best continuation is weak. We use $\theta=0.5$ for the main range-event rankings and export sensitivity summaries for $\theta\in\{0.3,0.4,0.5,0.6,0.7\}$.

The threshold has a different interpretation for range and domain scores. For range IoU, $\theta=0.5$ is an intuitive mid-scale cutoff: a link must have substantial footprint overlap under the common range-space frame. Domain overlap is more dataset-dependent because vertex-set sizes, mesh sampling, and sheet containment behavior differ. Therefore the viewer also computes domain thresholds from percentiles of best available domain continuations. A median threshold $Q_{50}$ means that a visible domain link is at least as strong as a typical best continuation in the current domain graph; higher percentiles such as $Q_{75}$ and $Q_{90}$ show increasingly conservative domain continuity.

\myparagraph{Continuing features.}
For each sheet, we greedily follow the best outgoing link while the score remains above $\theta$. The resulting lifetime is the number of timesteps visited. This produces a practical persistence diagnostic for prominent sheets, but it does not enforce global one-to-one identity.

\myparagraph{Domain/range target disagreement.}
For each source sheet, the viewer compares the target chosen by range shape with the target chosen by domain overlap. If the two targets differ, the source contributes a disagreement example. The score measures how much each mode loses when forced to accept the other mode's target, multiplied by a confidence term. The graph plots, for each timestep pair, the maximum source-level disagreement score. A high value means that at least one source sheet has strong but conflicting range and domain continuation evidence.

\section{Datasets}

The evaluation contains two real molecular datasets, MVK and stilbene, and one controlled torus dataset. MVK is represented by two generated artifacts, \texttt{MVK\_s1} and \texttt{MVK\_s2}, corresponding to the two MVK stages analyzed in the viewer. The torus dataset is a synthetic validation case with one controlled transition.

\begin{table}[t]
  \centering
  \caption{Datasets in the current generated artifacts. Each timestep keeps the top $N=20$ sheets by range-space area. Zero-vertex sheets are retained for range analysis but have zero domain-overlap evidence.}
  \label{tab:datasets}
  \begin{tabular}{lrrrr}
    \toprule
    Dataset & Timesteps & Sheets & Top $N$ & Zero-vertex sheets \\
    \midrule
    \texttt{MVK\_s1} & 83 & 1660 & 20 & 120 \\
    \texttt{MVK\_s2} & 82 & 1640 & 20 & 45 \\
    \texttt{stilbene} & 704 & 14080 & 20 & 21 \\
    \texttt{torus} & 100 & 2000 & 20 & 661 \\
    \bottomrule
  \end{tabular}
\end{table}

\begin{figure*}[t]
  \wideplaceholder{1.65in}{[Figure content: dataset overview. For MVK and stilbene, show a molecular rendering and representative sheet footprint/fiber-surface detail. For torus, show the synthetic volume and the controlled transition region. Include small labels with timestep counts and field names.]}
  \caption{Representative datasets used in the evaluation. MVK and stilbene are the two real molecular datasets, while the torus dataset provides one controlled transition case for validation.}
  \label{fig:datasets}
\end{figure*}

\section{Results}

We report range-space events as the primary feature-reconfiguration signal, because the current combined range score is shape IoU. Domain event rankings are reported separately as domain-overlap changes. Domain/range disagreement is used to select concrete examples where the two matching modes choose different target sheets.

\subsection{MVK: Recovering the Transition Window}
\label{sec:results-mvk}

The MVK trajectory was previously analyzed using continuous-scatterplot and image-moment descriptors~\cite{cspImageMomentsMVK}. That study reports a prominent transition near the approach to the conical intersection, with dense sampling around $29$--$31\,\mathrm{fs}$ and changes associated with the carbonyl oxygen and vinyl C3/C4 atoms. Our goal is not to reproduce atom tracking directly. Instead, we ask whether Reeb-space sheet tracking exposes the same transition as a feature-level event.

The viewer labels MVK timesteps by simulation step and converts them to femtoseconds using $t_\mathrm{fs}=\mathrm{step}/41.341374575751$. Thus the dense window $1240$--$1280$ corresponds to approximately $30.00$--$30.96\,\mathrm{fs}$.

\begin{table*}[t]
  \centering
  \caption{Important MVK range-space intervals identified by Reeb-space sheet tracking. Range events use shape IoU with $\theta=0.5$. The current \texttt{MVK\_s2} artifacts skip \texttt{step\_01360}, so the post-transition interval spans $1340$--$1380$.}
  \label{tab:mvk-range-events}
  \begin{tabular}{llrrrl}
    \toprule
    Dataset & Mode & Interval & Time (fs) & Score & Observation \\
    \midrule
    \texttt{MVK\_s1} & range & 1200--1220 & 29.03--29.51 & 35.77 & strongest \texttt{MVK\_s1} continuation loss \\
    \texttt{MVK\_s1} & range & 1180--1200 & 28.54--29.03 & 34.63 & onset of strong range reconfiguration \\
    \texttt{MVK\_s1} & range & 1276--1280 & 30.86--30.96 & 31.70 & late dense-window reconfiguration \\
    \texttt{MVK\_s2} & range & 1140--1160 & 27.58--28.06 & 37.21 & earliest strong \texttt{MVK\_s2} range reconfiguration \\
    \texttt{MVK\_s2} & range & 1340--1380 & 32.41--33.38 & 33.14 & post-transition reconfiguration \\
    \texttt{MVK\_s2} & range & 1160--1180 & 28.06--28.54 & 30.71 & continuation of the main event window \\
    \bottomrule
  \end{tabular}
\end{table*}

The strongest \texttt{MVK\_s1} range events occur at $29.03$--$29.51\,\mathrm{fs}$, $28.54$--$29.03\,\mathrm{fs}$, and $30.86$--$30.96\,\mathrm{fs}$. \texttt{MVK\_s2} shows strong range events at $27.58$--$28.06\,\mathrm{fs}$ and $28.06$--$28.54\,\mathrm{fs}$, with a later post-transition event at $32.41$--$33.38\,\mathrm{fs}$. This agreement with the known transition window suggests that range-space Reeb sheets capture a feature-level reconfiguration of the bivariate molecular field.

\begin{figure*}[t]
  \wideplaceholder{1.8in}{[Figure content: range-mode Sankey views for \texttt{MVK\_s1} and \texttt{MVK\_s2} over $27.5$--$31.5\,\mathrm{fs}$. Show the event-score graph above the Sankey panel. Highlight \texttt{MVK\_s1} intervals 1200--1220, 1180--1200, and 1276--1280, and \texttt{MVK\_s2} intervals 1140--1160 and 1160--1180.]}
  \caption{Range-space Reeb-sheet events concentrate near the MVK transition window reported by prior continuous-scatterplot analysis.}
  \label{fig:mvk-range-events}
\end{figure*}

The range metric also identifies persistent feature families. In \texttt{MVK\_s1}, three high-ranking sheets persist across the full time span $0.00$--$35.80\,\mathrm{fs}$: sheet $0\rightarrow6341$ with mean continuation score $0.978$, sheet $5251\rightarrow8$ with mean score $0.950$, and sheet $4\rightarrow5$ with mean score $0.896$. In \texttt{MVK\_s2}, sheet $24\rightarrow503$ persists across the full trajectory with mean score $0.900$. Two other prominent \texttt{MVK\_s2} tracks are sheet $22\rightarrow1892$, which persists until $27.58\,\mathrm{fs}$, and sheet $2\rightarrow0$, which starts at $13.06\,\mathrm{fs}$ and persists to the end.

The stride check supports the same conclusion. For \texttt{MVK\_s1}, the top range intervals expand from $1200$--$1220$ at stride 1 to $1180$--$1220$, $1160$--$1220$, and $1140$--$1220$ at strides 2, 3, and 4. For \texttt{MVK\_s2}, the top intervals expand from $1140$--$1160$ to $1140$--$1180$, $1120$--$1180$, and $1100$--$1180$. Coarser strides therefore cover the same temporal transition instead of moving the event elsewhere.

Domain-space tracking answers a different question. Using the maximum directional domain-overlap ratio, the strongest MVK domain events identify intervals where the same regular domain vertices stop overlapping the same sheet behavior: \texttt{MVK\_s1} has strong domain events at $500$--$520$, $860$--$880$, and $820$--$840$, while \texttt{MVK\_s2} has strong domain events at $960$--$980$, $520$--$540$, and $540$--$560$. These do not duplicate the range-transition timing. We interpret them as a separate domain-overlap result that should be inspected with linked sheet and fiber-surface views before assigning atom-specific meaning.

Domain/range target disagreement provides targeted examples near the transition. Strong disagreement examples include \texttt{MVK\_s1} intervals $1240$--$1244$ and $1220$--$1240$, and \texttt{MVK\_s2} intervals $1244$--$1248$, $1220$--$1240$, and $1248$--$1252$. These examples indicate that during the transition, the range-best target and domain-best target can differ for the same source sheet. The linked view is therefore useful for selecting concrete sheet pairs to inspect in the molecular context.

\begin{figure*}[t]
  \wideplaceholder{1.65in}{[Figure content: linked sheet/fiber-surface detail views for MVK transition examples. Suggested examples: \texttt{MVK\_s1} interval 1240--1244 with source sheet 978, and \texttt{MVK\_s2} interval 1244--1248 with source sheet 144. Show source sheet, range-selected target, domain-selected target, and the corresponding fiber-surface images.]}
  \caption{Domain/range disagreement examples near the MVK transition. The figure should show where range and domain criteria choose different continuations for the same source sheet.}
  \label{fig:mvk-disagreement-details}
\end{figure*}

\subsection{Stilbene: Long Molecular Sheet Evolution}
\label{sec:results-stilbene}

The stilbene dataset provides a longer molecular trajectory for testing whether Reeb-space sheet tracking remains useful beyond the MVK transition case. We again use the top $20$ sheets and the range-space shape-IoU score.

\begin{table}[t]
  \centering
  \caption{Important stilbene range-space intervals with $\theta=0.5$.}
  \label{tab:stilbene-events}
  \begin{tabular}{lrrrl}
    \toprule
    Interval & Score & Mean best & Weak src/tgt & Observation \\
    \midrule
    11915--11920 & 54.49 & 0.276 & 20/20 & strongest range reconfiguration \\
    11920--11925 & 53.44 & 0.328 & 20/20 & continuation of late event \\
    9380--9400 & 35.94 & 0.478 & 13/12 & earlier secondary event \\
    9660--9680 & 28.99 & 0.550 & 9/9 & secondary local reconfiguration \\
    11910--11915 & 28.13 & 0.494 & 9/9 & onset of late event \\
    2420--2440 & 27.65 & 0.518 & 9/9 & earlier secondary event \\
    \bottomrule
  \end{tabular}
\end{table}

The dominant stilbene range-space event occurs at $11915$--$11920$, followed by $11920$--$11925$. Unlike MVK, we do not currently tie these intervals to an external chemical transition label. Instead, they demonstrate that the tracking pipeline can isolate sheet reconfiguration intervals in a long molecular sequence.

The late event is robust. It remains the top event for thresholds $\theta=0.3$ through $\theta=0.7$. The stride check also preserves the same late event: the top intervals at strides 1, 2, 3, and 4 are $11915$--$11920$, $11910$--$11920$, $11905$--$11920$, and $11900$--$11920$. Thus coarser sampling expands the interval around the same endpoint rather than producing an unrelated peak.

The event intervals coexist with long-lived sheet tracks. The longest threshold-$0.5$ tracks include sheet $86\rightarrow6270$ over labels $2100$--$5340$, with length 163 and mean continuation score $0.860$; sheet $38\rightarrow19239$ over labels $9400$--$11255$, with length 162 and mean score $0.844$; and sheet $19\rightarrow53844$ over labels $5540$--$8700$, with length 159 and mean score $0.916$. These tracks show that prominent range-space sheets can remain coherent over long portions of the trajectory even when short intervals exhibit strong reconfiguration.

The domain event ranking identifies a separate set of intervals where overlap among regular domain vertices changes strongly. Using the maximum directional domain-overlap ratio, the strongest domain events are $8440$--$8460$, $8420$--$8440$, $8520$--$8540$, $8500$--$8520$, and $8460$--$8480$. These intervals do not duplicate the range-event ranking. Domain/range target disagreement is useful for selecting concrete visual checks, with strong examples at $8940$--$8960$, $9060$--$9080$, and $5260$--$5280$.

\begin{figure*}[t]
  \wideplaceholder{1.8in}{[Figure content: stilbene range-event overview. Include one panel around 9380--9680 highlighting 9380--9400 and 9660--9680, and another around 11900--11935 highlighting 11915--11920 and 11920--11925. Show the event graph and corresponding Sankey columns.]}
  \caption{Stilbene range-space event intervals identify a robust late sheet reconfiguration and earlier secondary events.}
  \label{fig:stilbene-range-events}
\end{figure*}

\begin{figure}[t]
  \placeholder{1.45in}{[Figure content: stilbene long-lived tracks. Highlight tracks $86\rightarrow6270$, $38\rightarrow19239$, and $19\rightarrow53844$, and include one clicked sheet/fiber detail view from the late event interval.]}
  \caption{Long-lived stilbene Reeb-sheet tracks coexist with strong range-event intervals.}
  \label{fig:stilbene-long-lived}
\end{figure}

\subsection{Synthetic Torus Validation}
\label{sec:results-torus}

The torus dataset provides a controlled validation case. Unlike the two real molecular datasets, its purpose is not chemical interpretation but checking whether the event diagnostic localizes a known sheet reconfiguration.

\begin{table}[t]
  \centering
  \caption{Top torus range-space event intervals with $\theta=0.5$. The strongest scores are centered around the controlled transition near timestep 50.}
  \label{tab:torus-events}
  \begin{tabular}{lrrrl}
    \toprule
    Interval & Score & Mean best & Weak src/tgt & Observation \\
    \midrule
    50--51 & 49.08 & 0.246 & 15/15 & strongest transition interval \\
    49--50 & 49.05 & 0.248 & 15/15 & symmetric pre-transition interval \\
    51--52 & 48.09 & 0.296 & 15/15 & symmetric post-transition interval \\
    48--49 & 48.09 & 0.296 & 15/15 & pre-transition shoulder \\
    47--48 & 45.57 & 0.322 & 13/13 & broader transition region \\
    52--53 & 45.57 & 0.322 & 13/13 & broader transition region \\
    \bottomrule
  \end{tabular}
\end{table}

The strongest range-space event occurs at $50$--$51$, with nearly identical scores on the adjacent intervals $49$--$50$, $51$--$52$, and $48$--$49$. This symmetric event cluster is the expected behavior for the synthetic sequence and provides a sanity check for the event diagnostic.

The torus event is stable for all tested thresholds from $\theta=0.3$ through $\theta=0.7$: the top interval remains $50$--$51$. The stride views preserve the same transition center. At strides 1, 2, 3, and 4, the top range intervals are $50$--$51$, $50$--$52$, $50$--$53$, and $50$--$54$. Coarser sampling therefore expands around the controlled transition instead of shifting the detected event.

The tracker also identifies long-lived torus sheet families. Several threshold-$0.5$ tracks span the full sequence, including $379\rightarrow43$, $383\rightarrow43$, $1183\rightarrow43$, and $13903\rightarrow43$, each with length 100 and mean continuation score approximately $0.926$. These tracks validate temporal continuity but also illustrate the non-one-to-one matching model: several tracks can converge to the same final sheet.

Domain events again answer a different question. Using the maximum directional domain-overlap ratio, the strongest domain events are $7$--$8$, $90$--$91$, $9$--$10$, $8$--$9$, and $92$--$93$, which indicate domain-vertex overlap behavior rather than the range transition at timestep 50. However, domain/range target disagreement peaks near the torus transition, including $48$--$49$ and $49$--$50$. This explains that range and domain criteria diverge near the controlled reconfiguration even though the range-event curve remains the cleanest validation signal.

\begin{figure*}[t]
  \wideplaceholder{1.75in}{[Figure content: torus validation figure. Show the event-score curve over all timesteps, highlight 45--55, mark 49--50, 50--51, and 51--52, and annotate stride 1--4 intervals centered at timestep 50.]}
  \caption{The torus event diagnostic localizes the controlled transition near timestep 50 and remains stable under threshold and stride changes.}
  \label{fig:torus-event-validation}
\end{figure*}

\section{Discussion}

\myparagraph{Range events are the primary feature-reconfiguration signal.}
Across the three case studies, range-space shape IoU gives the clearest feature-level event ranking. It recovers the known MVK transition window, identifies a robust late event in stilbene, and localizes the controlled torus transition. This supports using range-space Reeb-sheet geometry as the primary correspondence cue when the scientific question concerns persistence or reconfiguration of bivariate field relationships.

\myparagraph{Domain events are useful, but answer a different question.}
Domain-event rankings are not expected to match range-event rankings. They identify intervals where overlap among regular domain vertices changes strongly. This is useful for a different class of questions: whether the same vertices in the domain continue to participate in the sheet. We therefore keep domain event ranking as a separate analysis direction rather than using it as a validation or replacement for range-space tracking.

\myparagraph{Domain/range disagreement is an example-selection tool.}
The disagreement graph is most useful when it points to concrete source sheets for visual inspection. A high maximum disagreement score means that range and domain criteria confidently choose different targets for at least one source sheet in the interval. Such cases help explain why the two views diverge. Aggregate correlation summaries were removed from the current analysis because they did not provide a clear visual or scientific interpretation.

\myparagraph{Zero-vertex sheets.}
Some top-ranked sheets have no stored regular vertices. We retain them because they can have meaningful range-space geometry, and removing them would bias the range-space sheet ranking. However, they provide no evidence for vertex-overlap continuity. Domain conclusions should therefore be interpreted as regular-vertex overlap evidence, not as a complete measure of every possible sheet preimage contribution.

\section{Limitations}

The method depends on successful Reeb-space computation and geometry export. Degenerate input can cause missing timesteps or missing sheet geometry, and such timesteps are excluded from the temporal graph. The correspondence scores are heuristic and adjacent-pair based; the method does not solve a global assignment over the full sequence. Greedy lifetimes can therefore fragment or converge when multiple sheets have similar outgoing scores. The event score is a visual analytics diagnostic, not a formal topological event detector. Finally, domain overlap is currently based on stored regular vertices from the RSI representation, which limits precision near singular structures and for zero-vertex sheets.

\section{Conclusion}

We presented a Reeb-space sheet tracking workflow for time-varying bivariate fields. The method treats prominent Reeb-space sheets as feature units, computes range-space and domain-space candidate correspondences, and exposes the resulting temporal graph through a linked visual analytics interface. The results show that range-space Reeb-sheet events can recover known or controlled reconfiguration intervals, while domain-space rankings and domain/range disagreement provide complementary ways to inspect domain-vertex overlap changes. This supports Reeb-space sheets as practical trackable features for visual analysis of time-varying bivariate data.

\acknowledgments{}

\bibliographystyle{abbrv-doi}
\bibliography{bib}

\end{document}
